% !chktex-file 8

% LTeX: enabled=false

\documentclass[14pt,aspectratio=169,table,french]{beamer}
\usepackage{bxdpx-beamer}
\usepackage[french]{babel}
\usepackage{caption}
\usepackage{hyperref}
\usepackage{bookmark}
\usepackage{luatexja}
\usepackage{luatexja-fontspec}
\usepackage{fontspec}
\usepackage[T1]{fontenc}
\usepackage{lmodern}
\usepackage{newtxtext}
\usepackage{amsmath}
\usepackage{breqn}
\usepackage{silence}
\usepackage{stmaryrd}
\usepackage{hyperref}
\usepackage{tikz}
\usepackage{xcolor}
\usepackage{minted}
\usepackage{xcolor}
\usepackage{csquotes}

\usetikzlibrary{arrows, arrows.meta}
\usetikzlibrary{automata}
\usetikzlibrary{positioning}
\usetikzlibrary{shapes.geometric, shapes.misc}
\usetikzlibrary{math}
\usetikzlibrary{calc}
\usetikzlibrary{fit}

\WarningFilter{latexfont}{Font shape}

\newcommand{\myframetitle}{}
\newcommand{\setframetitle}[1]{\renewcommand{\myframetitle}{#1}}

\tikzset{
  ->,
  >=stealth',
  node distance=3cm,
  every state/.style={
    thick,
    rounded rectangle
  },
  initial text=\( \),
}

% Light gray background for code blocks
\definecolor{bgcolor}{HTML}{F0F0F0}

% Gray for function in code blocks
\definecolor{mgray}{HTML}{666666}

% Minted configuration
\setminted{
  linenos,              % Enable line numbers
  numbersep=8pt,        % Space between line numbers and code
  xleftmargin=2em,      % Add margin so numbers stay inside the block
  bgcolor=bgcolor,      % Set the background color of the code block
  style=algol_nu,       % Disable syntax highlighting (use plain black text)
  frame=none,           % Add a frame around the code block
  framesep=5pt,         % Space between the frame and the code
  breaklines=true,      % Automatically break long lines
  escapeinside=||       % Escape to LaTeX between the two characters specified
}

\DeclareMathOperator*{\NFA}{NFA}
\DeclareMathOperator*{\DFA}{DFA}
\DeclareMathOperator*{\RAT}{Rat}
\DeclareMathOperator*{\ndet}{ndet}
\DeclareMathOperator*{\nsdet}{(n)det} % chktex 36

\setbeamertemplate{theorems}[numbered]

\usetheme{Boadilla}
\usefonttheme{professionalfonts}

\DeclareFontFamily{OMX}{mlmex}{}
\DeclareFontShape{OMX}{mlmex}{m}{n}{<->mlmex10}{}
\renewcommand{\familydefault}{\sfdefault}
\renewcommand{\kanjifamilydefault}{\gtdefault}

\usefonttheme{structurebold}

\setbeamerfont{alerted text}{series=\bfseries}
\setbeamerfont{section in toc}{series=\mdseries}
\setbeamerfont{frametitle}{size=\Large}
\setbeamerfont{title}{size=\LARGE}
\setbeamerfont{date}{size=\small}

\definecolor{UniBlue}{RGB}{0,150,200}
\definecolor{AlertOrange}{RGB}{255,76,0}
\definecolor{AlmostBlack}{RGB}{38,38,38}
\setbeamercolor{normal text}{fg=AlmostBlack}
\setbeamercolor{structure}{fg=UniBlue}
\setbeamercolor{block title}{fg=UniBlue!50!black}
\setbeamercolor{alerted text}{fg=AlertOrange}
\mode<beamer>{
  \definecolor{BackGroundGray}{RGB}{254,254,254}
  \setbeamercolor{background canvas}{bg=BackGroundGray}
}

\setbeamertemplate{blocks}[rounded]
\useinnertheme{circles}
\setbeamertemplate{navigation symbols}{}
\setbeamertemplate{footline}[frame number]
\setbeamertemplate{frametitle continuation}{}

\setbeamertemplate{title page}{
  \vspace{2.0em}
  {\usebeamerfont{title} \usebeamercolor[fg]{title} \inserttitle~\par}
  {\usebeamerfont{subtitle}\usebeamercolor[fg]{subtitle}\insertsubtitle~\par}
  \vspace{2.5em}
  \begin{minipage}{0.48\textwidth}
    \usebeamercolor[fg]{titlegraphic}\inserttitlegraphic{}
  \end{minipage}
  \hfill
  \begin{minipage}{0.48\textwidth}
    \begin{flushright}
      \usebeamerfont{author}\insertauthor\par
      \usebeamerfont{institute}\insertinstitute~\par
      \vspace{3em}
      \usebeamerfont{date}\insertdate\par
    \end{flushright}
  \end{minipage}
}

\newcounter{firstsection}
\setcounter{firstsection}{1}

\AtBeginSection[]{
  \ifnum\value{firstsection}=1
  \setcounter{firstsection}{0}
  \begin{frame}
    \tableofcontents[hideallsubsections]
  \end{frame}
  \else
  \begin{frame}
    \centering
    \usebeamerfont{title}
    \usebeamercolor[fg]{title}
    \insertsection{}
  \end{frame}
  \fi
}

% LTeX: enabled=true

\title{Complexité en état opérationnel de langage rationnel}
\subtitle{
  Master 1 - ITA~: Application Informatique \\ Encadré par Pascal Caron
}
\author{Edouard HADDAG}
\institute{Université de Rouen Normandie}
\date{\today}
\titlegraphic{
  \includegraphics[scale=0.3]{
    resources/logo-univ-rouen-normandie-noir.png
  }
}

\begin{document}

\maketitle

% !chktex-file 8

\subsection{Présentation du papier}

\setframetitle{Présentation du papier}

\begin{frame}{\myframetitle}
  \begin{block}{Présentation du papier}
    % LTex: language=en
    \textit{State complexity of some operations on binary regular
    languages}\footnotemark{}
    % LTex: language=fr

    \vphantom{}

    Par \textbf{Galina Jirásková}, publié dans
    % LTex: language=en
    \textit{\underline{Theoretical Computer Science, 330.2}}
    % LTex: language=fr%
    en 2005
  \end{block}
  \footnotetext{\hyperlink{JIRASKOVA2005287}{\cite{JIRASKOVA2005287}}}
\end{frame}

\section{Introduction}

\subsection{Prémisse}

\subsubsection{Alphabet et Mot}

\setframetitle{Alphabet et Mot}

\begin{frame}{\myframetitle}
  \begin{definition}[Alphabet]
    Un << alphabet >> \(\Sigma\) est un ensemble fini non-vide de symboles.
  \end{definition}

  \pause[]

  \begin{definition}[Mot]
    Un << mot >> \(w\) est une suite finie de symboles sur un alphabet
    \(\Sigma\). La suite vide de symbole sera notée par \(\varepsilon\).
  \end{definition}

  \pause[]

  \begin{example}
    \centering
    \vspace{-1.5\topsep}
    \begin{gather*}
      \Sigma = \{a, b\} \\
      w = abba
    \end{gather*}
  \end{example}
\end{frame}

\begin{frame}{\myframetitle}
  \begin{definition}[Concaténation d'un mot]
    L'opération \(\cdot\) désigne la << concaténation >> de deux mots.
  \end{definition}

  \pause[]

  \begin{example}
    \centering
    \vspace{-1.5\topsep}
    \begin{gather*}
      u = aa \text{ et } v = bb \\
      w = u \cdot v = aabb
    \end{gather*}
  \end{example}
\end{frame}

\begin{frame}{\myframetitle}
  \begin{definition}[Le miroir d'un mot]
    On notera \(\overleftarrow w\), le << miroir >> du mot \(w\).
  \end{definition}

  \pause[]

  \begin{example}
    \vspace{-1.5\topsep}
    \begin{align*}
      u &= abcd\\
      \overleftarrow u &= dcba
    \end{align*}
  \end{example}
\end{frame}

\subsubsection{Langage}

\setframetitle{Langage}

\begin{frame}{\myframetitle}
  \begin{definition}[Langage]
    Un << langage >> \(L\) est un ensemble de mots sur un alphabet \(\Sigma\).
  \end{definition}

  \pause[]

  \begin{example}
    \vspace{-1.5\topsep}
    \begin{gather*}
      L = \{\varepsilon, aa, bb, abab\}
    \end{gather*}
  \end{example}
\end{frame}

\begin{frame}{\myframetitle}
  \begin{definition}[La concaténation]
    La << concaténation >> est définie grâce à la concaténation des mots~:

    \begin{alignat*}{2}
      &L_1 &&= \{a, b\} \\
      &L_2 &&= \{c, d\} \\
      &L_1 \cdot L_2 &&= \{ac, ad, bc, bd\} \\
    \end{alignat*}
  \end{definition}
\end{frame}

\begin{frame}{\myframetitle}
  \begin{definition}[La copie n-ième]
    On définit la << copie n-ième >> d'un langage \(L\) noté \(L^n\)~:

    \begin{alignat*}{2}
      &L_1 &&= \{a, b\} \\
      &L_1^2 &&= \{aa, ab, bb, ba\}
    \end{alignat*}
  \end{definition}
\end{frame}

\begin{frame}{\myframetitle}
  \begin{definition}[L'étoile (de \textit{Kleene}) d'un langage]
    On peut définir << l'étoile >> d'un langage notée \(L^*\)~:
    % LTeX: enabled=false
    \begin{align*}
      L^* = \bigcup_{i \geq 0} L^i
    \end{align*}
    % LTeX: enabled=true
  \end{definition}

  \pause[]

  \begin{example}
    \vspace{-1.5\topsep}
    \begin{align*}
      &L = \{a, b\} \\
      &L^* = \{\varepsilon, a, b, aa, ab, \ldots\}
    \end{align*}
  \end{example}
\end{frame}

\begin{frame}{\myframetitle}
  \begin{definition}[Le langage de l'ensemble des mots]
    L'ensemble des mots possible sur l'alphabet \(\Sigma\) sera noté
    \(\Sigma^*\).
  \end{definition}

  \pause[]

  \begin{definition}[Le langage miroir]
    On pourra noter \(\overleftarrow L\), << le langage miroir >> de \(L\), qui
    sera le langage des miroirs des mots de \(L\)~:
    \begin{align*}
      \overleftarrow L = \{\overleftarrow w ~|~ w \in L\}
    \end{align*}
  \end{definition}

  \pause[]

  \begin{block}{Autres opérations}
    Les langages étant des ensembles, leurs opérations peuvent être étendues
    aux langages.
  \end{block}
\end{frame}

\begin{frame}{\myframetitle}
  \begin{definition}[Les langages rationnels]
    L'ensemble des << langages rationnels >> sur \(\Sigma\), noté
    \(\RAT(\Sigma^*)\), est le plus petit ensemble de langages sur \(\Sigma\)
    qui vérifie~:

    \vphantom{}

    \begin{itemize}
      \item \(\RAT(\Sigma^*)\) contient \(\varnothing\) et \(\{a\}\) avec \(a
        \in \Sigma\).
      \item \(\RAT(\Sigma^*)\) est fermé pour l'union, la concaténation et
        l'étoile.
    \end{itemize}
  \end{definition}
\end{frame}

\subsubsection{Automate}

\setframetitle{Automate}

\begin{frame}{\myframetitle}
  \begin{definition}[Automate]
    \(M = (Q, I, F, \delta)\) avec \(Q = \{q_1, q_2, q_3\}\), \(I = \{q_1\}\) et
    \(F = \{q_3\}\).
    \begin{center}
      \begin{tikzpicture}
        \node[state, initial] (q1) {\(q_1\)};
        \node[state, right of=q1] (q2) {\(q_2\)};
        \node[state, accepting, right of=q2] (q3) {\(q_3\)};

        \draw
        (q1) edge[loop above] node{\(a\)} (q1)
        (q1) edge[above] node{\(b\)} (q2)
        (q1) edge[bend right, below] node{\(a\)} (q3)
        (q2) edge[above] node{\(a,b\)} (q3)
        (q3) edge[bend right=50, above] node{\(b\)} (q2);
      \end{tikzpicture}
    \end{center}
  \end{definition}
\end{frame}

\begin{frame}{\myframetitle}
  \begin{definition}[L'acceptation d'un mot]
    Le mot \(aba\) est << accepté >>~:
    \begin{center}
      \begin{tikzpicture}
        \node[state, initial] (q1) {\(q_1\)};
        \node[state, right of=q1] (q2) {\(q_2\)};
        \node[state, accepting, right of=q2] (q3) {\(q_3\)};

        \only<1> {\draw (q1) edge[loop above] node{\(a\)} (q1);}
        \only<1-2> {\draw (q1) edge[above] node{\(b\)} (q2);}
        \only<1-3> {\draw (q2) edge[above] node{\(a,b\)} (q3);}

        \only<2-> {\draw (q1) edge[loop above, red] node{\(a\)} (q1);}
        \only<3-> {\draw (q1) edge[above, red] node{\(b\)} (q2);}
        \only<4-> {\draw (q2) edge[above, red] node{\(a,b\)} (q3);}

        \draw
        (q3) edge[bend right=50, above] node{\(b\)} (q2)
        (q1) edge[bend right, below] node{\(a\)} (q3);
      \end{tikzpicture}
    \end{center}
  \end{definition}
\end{frame}

\begin{frame}{\myframetitle}
  \begin{definition}[Automate déterministe]
    On parlera d'automate << déterministe >> quand~:
    \begin{center}
      \only<1> {
        \begin{tikzpicture}
          \node[state, initial] (q1) {\(q_1\)};
          \node[state, right of=q1] (q2) {\(q_2\)};
          \node[state, accepting, right of=q2] (q3) {\(q_3\)};

          \draw
          (q1) edge[loop above] node{\(a\)} (q1)
          (q1) edge[above] node{\(b\)} (q2)
          (q2) edge[above] node{\(a,b\)} (q3)
          (q3) edge[bend right=50, above] node{\(b\)} (q2)
          (q3) edge[bend left, below] node{\(a\)} (q1);
        \end{tikzpicture}
      }

      \only<2> {
        \begin{tikzpicture}[every initial by arrow/.style={red}]
          \node[state, initial] (q1) {\(q_1\)};
          \node[state, right of=q1] (q2) {\(q_2\)};
          \node[state, accepting, right of=q2] (q3) {\(q_3\)};

          \draw
          (q1) edge[loop above] node{\(a\)} (q1)
          (q1) edge[above] node{\(b\)} (q2)
          (q2) edge[above] node{\(a,b\)} (q3)
          (q3) edge[bend right=50, above] node{\(b\)} (q2)
          (q3) edge[bend left, below] node{\(a\)} (q1);
        \end{tikzpicture}
      }

      \only<3> {
        \begin{tikzpicture}[every initial by arrow/.style={gray}]
          \node[state, initial] (q1) {\(q_1\)};
          \node[state, right of=q1, gray] (q2) {\(q_2\)};
          \node[state, accepting, right of=q2, gray] (q3) {\(q_3\)};

          \draw
          (q1) edge[loop above] node{\(a\)} (q1)
          (q1) edge[above] node{\(b\)} (q2)
          (q2) edge[above, gray] node{\(a,b\)} (q3)
          (q3) edge[bend right=50, above, gray] node{\(b\)} (q2)
          (q3) edge[bend left, below, gray] node{\(a\)} (q1);
        \end{tikzpicture}
      }

      \only<4> {
        \begin{tikzpicture}[every initial by arrow/.style={gray}]
          \node[state, initial, gray] (q1) {\(q_1\)};
          \node[state, right of=q1] (q2) {\(q_2\)};
          \node[state, accepting, right of=q2, gray] (q3) {\(q_3\)};

          \draw
          (q1) edge[loop above, gray] node{\(a\)} (q1)
          (q1) edge[above, gray] node{\(b\)} (q2)
          (q2) edge[above] node{\(a,b\)} (q3)
          (q3) edge[bend right=50, above, gray] node{\(b\)} (q2)
          (q3) edge[bend left, below, gray] node{\(a\)} (q1);
        \end{tikzpicture}
      }

      \only<5> {
        \begin{tikzpicture}[every initial by arrow/.style={gray}]
          \node[state, initial, gray] (q1) {\(q_1\)};
          \node[state, right of=q1, gray] (q2) {\(q_2\)};
          \node[state, accepting, right of=q2] (q3) {\(q_3\)};

          \draw
          (q1) edge[loop above, gray] node{\(a\)} (q1)
          (q1) edge[above, gray] node{\(b\)} (q2)
          (q2) edge[above, gray] node{\(a,b\)} (q3)
          (q3) edge[bend right=50, above] node{\(b\)} (q2)
          (q3) edge[bend left, below] node{\(a\)} (q1);
        \end{tikzpicture}
      }
    \end{center}
  \end{definition}
\end{frame}

\begin{frame}{\myframetitle}
  \begin{definition}[Automate minimal]
    L'automate \(M_1\) est << minimal >> du langage \(L = \{ab\}\)~:

    \begin{minipage}{0.47\textwidth}
      \centering
      \begin{tikzpicture}[node distance=2.5cm]
        \node[state, initial] (0) {\(0\)};
        \node[state, right of=0] (1) {\(1\)};
        \node[state, accepting, below of=1] (2) {\(2\)};

        \node[above =.75cm of $(0)!0.5!(1)$, align=center] {\(M_1\)};

        \draw
        (0) edge[above] node{\(a\)} (1)
        (1) edge[right] node{\(b\)} (2);
      \end{tikzpicture}
    \end{minipage}
    \hfill
    \begin{minipage}{0.47\textwidth}
      \centering
      \begin{tikzpicture}[node distance=2.5cm]
        \node[state, initial] (0) {\(0\)};
        \node[state, right of=0] (1) {\(1\)};
        \node[state, accepting, below of=1] (2) {\(2\)};
        \node[state, left of=2] (3) {\(3\)};

        \node[above =.75cm of $(0)!0.5!(1)$, align=center] {\(M_2\)};

        \draw
        (0) edge[above] node{\(a\)} (1)
        (1) edge[right] node{\(b\)} (2)
        (0) edge[left] node{\(a\)} (3);
      \end{tikzpicture}
    \end{minipage}
  \end{definition}
\end{frame}

\subsubsection{Lien entre les langages et les automates}

\setframetitle{Lien entre les langages et les automates}

\begin{frame}{\myframetitle}
  \begin{definition}[L'ensemble des langages reconnaisable]
    On définira l'ensemble des langages de \(\Sigma^*\) << reconnaissable >> par
    au moins un automate par \(\REC(\Sigma^*)\).
  \end{definition}

  \pause[]

  \begin{block}{Théorème de \textit{Kleene}}
    L'informaticien \textit{Stephen C. Kleene} a montré en 1956 que~:
    % LTeX: enabled=false
    \begin{align*}
      \REC(\Sigma^*) = \RAT(\Sigma^*)
    \end{align*}
    % LTeX: enabled=true
  \end{block}
\end{frame}

\subsubsection{Complexité en état}

\setframetitle{Complexité en état}

\begin{frame}{\myframetitle}
  \begin{definition}
    La << complexité en état >> est une mesure d'un \textbf{langage}, elle est
    définie comme le nombre d'états de l'automate minimal du langage~:

    \vphantom{}

    \begin{itemize}
      \item Complexité en état déterministe (qu'on notera \(C_{\det}\)).
      \item Complexité en état non déterministe (qu'on notera \(C_{\ndet}\)).
    \end{itemize}
  \end{definition}
\end{frame}

\begin{frame}{\myframetitle}
  \begin{example}[%
      \only<1>{Définition d'une famille}%
      \only<2>{\(C_{\ndet}(A_2) = 2 + 2\)}%
      \only<3>{\(C_{\det}(A_2) = 2^{2 + 1}\)}%
    ]
    \centering
    \vspace{-1.5\topsep}
    \begin{gather*}
      A_n = \Sigma^* \cdot a \cdot \Sigma^n
    \end{gather*}
    \vspace{-2.25\topsep}

    \onslide<2-> {
      \begin{tikzpicture}[node distance=2.5cm, scale=0.7,transform shape]
        \node[state, initial] (0') {\(0\)};
        \node[state, right of=0'] (1') {\(1\)};
        \node[state, below of=1'] (2') {\(2\)};
        \node[state, accepting, left of=2'] (3') {\(3\)};

        \draw
        (0') edge[above] node{\(a\)} (1')
        (1') edge[right] node{\(a,b\)} (2')
        (2') edge[above] node{\(a,b\)} (3')
        (0') edge[loop above] node{\(a,b\)} (0');

        \onslide<3> {
          \node[state, initial, right =5cm of 1'] (0) {\(0\)};
          \node[state, right of=0] (1) {\(1\)};
          \node[state, right of=1] (2) {\(2\)};
          \node[state, right of=2] (3) {\(3\)};
          \node[state, accepting, below of=3] (4) {\(4\)};
          \node[state, accepting, left of=4] (5) {\(5\)};
          \node[state, accepting, left of=5] (6) {\(6\)};
          \node[state, accepting, left of=6] (7) {\(7\)};

          \draw
          (0) edge[loop above] node{\(b\)} (0)
          (3) edge[loop above] node{\(a\)} (3)

          (0) edge[above] node{\(a\)} (1)
          (1) edge[above] node{\(a\)} (2)
          (2) edge[above] node{\(a\)} (3)
          (3) edge[right] node{\(b\)} (4)
          (4) edge[above] node{\(a\)} (5)
          (5) edge[above] node{\(b\)} (6)
          (6) edge[above] node{\(b\)} (7)

          (1) edge[left] node{\(b\)} (6)
          (2) edge[left] node{\(b\)} (4)
          (4) edge[below, bend left] node{\(b\)} (7)
          (5) edge[left] node{\(a\)} (2)
          (6) edge[above, bend right] node{\(a\)} (5)
          (7) edge[left] node{\(a\)} (1)
          (7) edge[left] node{\(b\)} (0);
        }
      \end{tikzpicture}
    }
  \end{example}
\end{frame}

\begin{frame}{\myframetitle}
  \begin{definition}[Complexité en état opérationnielle]
    La fonction \(f\) à une complexité \(g(n, m)\) si~:
    \begin{center}
      \begin{tikzpicture}[transform shape, scale=0.8]
        \node[draw, ellipse, minimum width=1.5cm] (A) at (-2,2) {\(A\)};
        \node[draw, ellipse, minimum width=1.5cm] (B) at (2,2) {\(B\)};
        \node[draw, rectangle, minimum width=2cm, minimum height=1cm]
        (F) at (0,0) {\(f\)};
        \node[draw, ellipse, minimum width=1.5cm] (C) at (0,-2) {\(C\)};

        \node[left =0.25cm of A] {\(n\)};
        \node[right =0.25cm of B] {\(m\)};
        \node[below =0.25cm of C] {\(k \leq g(n, m)\)};

        \coordinate (center) at ($(F.north west)!0.5!(F.north east)$);

        \draw (A) -- (-0.5,2) -- ([shift={(-0.5cm,0cm)}] center);
        \draw (B) -- (0.5,2) -- ([shift={(0.5cm,0cm)}] center);
        \draw (F) -- (C);
      \end{tikzpicture}
    \end{center}
  \end{definition}
\end{frame}
\section{Le trognon d'un langage}

\subsection{Définition}

\setframetitle{Définition}

\begin{frame}{\myframetitle}
  \begin{block}{Le trognon d'un mot}
    \centering
    \includegraphics[width=0.25\textwidth]{./resources/full_apple.png}

    \pause[]

    \vphantom{}

    \Large{\(abcdcba\)}
  \end{block}
\end{frame}

\begin{frame}{\myframetitle}
  \begin{block}{Le trognon d'un mot}
    \centering
    \includegraphics[width=0.25\textwidth]{./resources/mid_apple.png}

    \pause[]

    \vphantom{}

    \Large{\(\color{red}{ab} \color{black}{\cdot cdc \cdot} \color{red}{ba}\)}
  \end{block}
\end{frame}

\begin{frame}{\myframetitle}
  \begin{block}{Le trognon d'un mot}
    \centering
    \includegraphics[width=0.25\textwidth]{./resources/empty_apple.png}

    \pause[]

    \vphantom{}

    \Large{\(\color{red}{abc} \color{black}{\cdot d \cdot} \color{red}{cba}\)}
  \end{block}
\end{frame}

\begin{frame}{\myframetitle}
  \begin{block}{Le trognon d'un mot}
    \centering
    \includegraphics[width=0.25\textwidth]{./resources/full_apple.png}

    \pause[]

    \vphantom{}

    \Large{\(\color{red}{\varepsilon} \color{black}{\cdot abcdcba \cdot}
    \color{red}{\varepsilon}\)}
  \end{block}
\end{frame}

\begin{frame}{\myframetitle}
  \begin{definition}[Le trognon d'un mot]
    Le trognon d'un langage sera alors l'union des trognons des mots qu'il le
    compose.
  \end{definition}
\end{frame}

\subsection{Algorithme de grignotage}

\setframetitle{Algorithme de grignotage}

\begin{frame}{\myframetitle}
  \begin{example}[\(L(M) = a \cdot a^* \cdot c \cdot a^* \cdot a\)]
    \begin{center}
      \begin{tikzpicture}
        \node[state, initial] (1) {\(1\)};
        \node[state, right of=1] (2) {\(2\)};
        \node[state, right of=2] (3) {\(3\)};
        \node[state, accepting, right of=3] (4) {\(4\)};

        \draw
        (1) edge[above] node{\(a\)} (2)
        (2) edge[loop above] node{\(a\)} (2)
        (2) edge[above] node{\(c\)} (3)
        (3) edge[loop above] node{\(a\)} (3)
        (3) edge[above] node{\(a\)} (4);
      \end{tikzpicture}
    \end{center}
  \end{example}
\end{frame}

\begin{frame}{\myframetitle}
  \begin{example}[\(M\) grignoté de \(a\)]
    \begin{center}
      \begin{tikzpicture}
        \node[state] (1) {\(1\)};
        \node[state, initial below, right of=1] (2) {\(2\)};
        \node[state, accepting, right of=2] (3) {\(3\)};
        \node[state, right of=3] (4) {\(4\)};

        \draw
        (1) edge[above, red] node{\(a\)} (2)
        (2) edge[loop above] node{\(a\)} (2)
        (2) edge[above] node{\(c\)} (3)
        (3) edge[loop above] node{\(a\)} (3)
        (3) edge[above, red] node{\(a\)} (4);
      \end{tikzpicture}
    \end{center}
  \end{example}
\end{frame}

\begin{frame}{\myframetitle}
  \begin{example}[\(M\) grignoté de \(aa\)]
    \begin{center}
      \begin{tikzpicture}
        \node[state] (1) {\(1\)};
        \node[state, initial below, right of=1] (2) {\(2\)};
        \node[state, accepting, right of=2] (3) {\(3\)};
        \node[state, right of=3] (4) {\(4\)};

        \draw
        (1) edge[above, red] node{\(a\)} (2)
        (2) edge[loop above, red] node{\(a\)} (2)
        (2) edge[above] node{\(c\)} (3)
        (3) edge[loop above, red] node{\(a\)} (3)
        (3) edge[above, red] node{\(a\)} (4);
      \end{tikzpicture}
    \end{center}
  \end{example}
\end{frame}

\begin{frame}{\myframetitle}
  \begin{definition}
    L'automate \(\texttt{nibbling}(M)\) sera donc l'union des automates
    grignotés distincts.
  \end{definition}

  \pause[]

  \begin{block}{Compléxité en état de notre algorithme}
    Le nombre d'états de l'union de deux automates est égal à la somme des
    nombres d'états des automates.

    \pause[]

    \vphantom{}

    Les automates grignotés ont comme forme \((Q, I, F, \delta)\) avec \(Q\)
    et \(\delta\) les mêmes et comme changement \(I\) et \(F\) qui sont deux
    ensembles non vides.

    \pause[]

    \vphantom{}

    Ce qui nous donne comme complexité finale~: \(n(2^n - 1)^2\) états.
  \end{block}
\end{frame}


\section{Langage permuté}

\subsection{Définition}

\setframetitle{Définition}

\begin{frame}{\myframetitle}
  \begin{definition}[\(Twist(L)\)]
    Le langage \(Twist(L)\) sera défini comment l'ensemble des mots du
    langage de \(L\) tel qu'on aura échangé les lettres aux indices \(2k\) et
    \(2k + 1\) avec \(k\in \mathbb{N}\).
  \end{definition}

  \pause[]

  \begin{example}
    \vspace*{-\baselineskip}
    \begin{align*}
      L &= \{\varepsilon, a, ab, abcd\} \\
      Twist(L) &= \{\varepsilon, a, ba, badc\}
    \end{align*}
  \end{example}
\end{frame}

\subsection{Algorithme twister}

\setframetitle{Algorithme twister}

\begin{frame}{\myframetitle}
  \begin{example}[Soit l'automate \(M\) suivant]
    \centering
    \begin{tikzpicture}
      \tikzset{
        ->,
        >=stealth',
        node distance=2.5cm,
        every state/.style={thick},
        initial text=$ $,
      }
      \node[state, initial] (q1) {\(1\)};
      \node[state, right of=q1] (q2) {\(2\)};
      \node[state, accepting, right of=q2] (q3) {\(3\)};
      \node[state, above right of=q3] (q4) {\(4\)};
      \node[state, right of=q4] (q5) {\(5\)};
      \node[state, accepting, right of=q5] (q6) {\(6\)};
      \node[state, right of=q3] (q7) {\(7\)};
      \node[state, accepting, right of=q7] (q8) {\(8\)};
      \node[state, right of=q8] (q9) {\(9\)};
      \node[state, accepting, below right of=q3] (q10) {\(10\)};
      \node[state, right of=q10] (q11) {\(11\)};
      \node[state, right of=q11] (q12) {\(12\)};

      \only<1> {
        \draw
          (q1) edge[above] node[blue]{\(b\)} (q2)
          (q2) edge[above] node[red]{\(a\)} (q3)
          (q3) edge[above] node[blue]{\(b\)} (q4)
          (q4) edge[above] node[red]{\(a\)} (q5)
          (q5) edge[above] node[blue]{\(b\)} (q6)
          (q3) edge[above] node[blue]{\(a\)} (q7)
          (q7) edge[above] node[red]{\(b\)} (q8)
          (q8) edge[above] node[blue]{\(b\)} (q9)
          (q3) edge[above] node[blue]{\(b\)} (q10)
          (q10) edge[above] node[red]{\(b\)} (q11)
          (q11) edge[above] node[blue]{\(a\)} (q12);
      }
      \only<2> {
        \draw
          (q1) edge[above] node[red]{\(a\)} (q2)
          (q2) edge[above] node[blue]{\(b\)} (q3)
          (q3) edge[above] node[red]{\(a\)} (q4)
          (q4) edge[above] node[blue]{\(b\)} (q5)
          (q5) edge[above] node[blue]{\(b\)} (q6)
          (q3) edge[above] node[red]{\(b\)} (q7)
          (q7) edge[above] node[blue]{\(a\)} (q8)
          (q8) edge[above] node[blue]{\(b\)} (q9)
          (q3) edge[above] node[red]{\(b\)} (q10)
          (q10) edge[above] node[blue]{\(b\)} (q11)
          (q11) edge[above] node[blue]{\(a\)} (q12);
      }
    \end{tikzpicture}
  \end{example}
\end{frame}

\begin{frame}{\myframetitle}
  \begin{example}[Soit l'automate \(N\) suivant]
    \centering
    \only<1-2> {
      \begin{tikzpicture}
        \tikzset{
          ->,
          >=stealth',
          node distance=2.25cm,
          every state/.style={
            thick,
            inner sep = 0pt,
            minimum size = 0.75cm,
          },
          initial text=$ $,
        }
        \node[state, initial] (q1) {\(1\)};
        \node[state, right of=q1] (q2) {\(2\)};
        \node[state, above right of=q2] (q3) {\(3\)};
        \node[state, right of=q3] (q4) {\(4\)};
        \node[state, accepting, right of=q4] (q5) {\(5\)};
        \node[state, right of=q2] (q6) {\(6\)};
        \node[state, accepting, right of=q6] (q7) {\(7\)};
        \node[state, right of=q7] (q8) {\(8\)};

        \draw
          (q1) edge[above] node[blue]{\(b\)} (q2)
          (q2) edge[above] node[red]{\(a\)} (q3)
          (q3) edge[above] node[blue]{\(b\)} (q4)
          (q4) edge[above] node[red]{\(a\)} (q5)
          (q2) edge[above] node[red]{\(b\)} (q6)
          (q6) edge[above] node[blue]{\(a\)} (q7)
          (q7) edge[above] node[red]{\(b\)} (q8);
      \end{tikzpicture}

      \pause[]

      \vphantom{}

      \begin{tikzpicture}
        \tikzset{
          ->,
          >=stealth',
          node distance=2.25cm,
          every state/.style={
            thick,
            inner sep = 0pt,
            minimum size = 0.75cm,
          },
          initial text=$ $,
        }
        \node[state, initial] (q1) {\(1\)};
        \node[state, right of=q1] (q2) {\(2\)};
        \node[state, above right of=q2] (q3) {\(3\)};
        \node[state, right of=q3] (q4) {\(4\)};
        \node[state, accepting, right of=q4] (q5) {\(5\)};
        \node[state, right of=q2] (q6) {\(6\)};
        \node[state, accepting, right of=q6] (q7) {\(7\)};
        \node[state, right of=q7] (q8) {\(8\)};

        \draw
          (q1) edge[above] node[red]{\(a, b\)} (q2)
          (q2) edge[above] node[blue]{\(b\)} (q3)
          (q3) edge[above] node[red]{\(a\)} (q4)
          (q4) edge[above] node[blue]{\(b\)} (q5)
          (q2) edge[above] node[blue]{\(b\)} (q6)
          (q6) edge[above] node[red]{\(b\)} (q7)
          (q7) edge[above] node[blue]{\(a\)} (q8);
      \end{tikzpicture}
    }
    \only<3-4> {
      \begin{tikzpicture}
        \tikzset{
          ->,
          >=stealth',
          node distance=2.25cm,
          every state/.style={
            thick,
            inner sep = 0pt,
            minimum size = 0.75cm,
          },
          initial text=$ $,
        }
        \node[state, initial] (q1) {\(1\)};
        \node[state, right of=q1] (q2b) {\(2_b\)};
        \node[state, left of=q3] (q2a) {\(2_a\)};
        \node[state, above right of=q2b] (q3) {\(3\)};
        \node[state, right of=q3] (q4) {\(4\)};
        \node[state, accepting, right of=q4] (q5) {\(5\)};
        \node[state, right of=q2] (q6) {\(6\)};
        \node[state, accepting, right of=q6] (q7) {\(7\)};
        \node[state, right of=q7] (q8) {\(8\)};

        \draw
          (q1) edge[above] node[blue]{\(b\)} (q2a)
          (q1) edge[above] node[blue]{\(b\)} (q2b)
          (q2a) edge[above] node[red]{\(a\)} (q3)
          (q3) edge[above] node[blue]{\(b\)} (q4)
          (q4) edge[above] node[red]{\(a\)} (q5)
          (q2b) edge[above] node[red]{\(b\)} (q6)
          (q6) edge[above] node[blue]{\(a\)} (q7)
          (q7) edge[above] node[red]{\(b\)} (q8);
      \end{tikzpicture}

      \pause[4]

      \vphantom{}

      \begin{tikzpicture}
        \tikzset{
          ->,
          >=stealth',
          node distance=2.25cm,
          every state/.style={
            thick,
            inner sep = 0pt,
            minimum size = 0.75cm,
          },
          initial text=$ $,
        }
        \node[state, initial] (q1) {\(1\)};
        \node[state, right of=q1] (q2b) {\(2_b\)};
        \node[state, left of=q3] (q2a) {\(2_a\)};
        \node[state, above right of=q2b] (q3) {\(3\)};
        \node[state, right of=q3] (q4) {\(4\)};
        \node[state, accepting, right of=q4] (q5) {\(5\)};
        \node[state, right of=q2] (q6) {\(6\)};
        \node[state, accepting, right of=q6] (q7) {\(7\)};
        \node[state, right of=q7] (q8) {\(8\)};

        \draw
          (q1) edge[above] node[red]{\(a\)} (q2a)
          (q1) edge[above] node[red]{\(b\)} (q2b)
          (q2a) edge[above] node[blue]{\(b\)} (q3)
          (q3) edge[above] node[red]{\(a\)} (q4)
          (q4) edge[above] node[blue]{\(b\)} (q5)
          (q2b) edge[above] node[blue]{\(b\)} (q6)
          (q6) edge[above] node[red]{\(b\)} (q7)
          (q7) edge[above] node[blue]{\(a\)} (q8);
      \end{tikzpicture}
    }
  \end{example}
\end{frame}

\begin{frame}{\myframetitle}
  \begin{block}{Automate quelconque}
    Si on ajoute les cycles, on devra alors supposer que tous n\oe{}uds
    intermédiaire se trouve dans le cas du n\oe{}ud \(2\), on devra donc le
    dupliquer.
  \end{block}
\end{frame}

\begin{frame}{\myframetitle}
  \begin{example}
    \centering
    \only<1-2> {
      \begin{tikzpicture}
        \tikzset{
          ->,
          >=stealth',
          node distance=2.25cm,
          every state/.style={
            thick,
            inner sep = 0pt,
            minimum size = 0.75cm,
          },
          initial text=$ $,
        }
        \node[state, initial] (1) {\(1\)};
        \node[state, right of=1] (2) {\(2\)};
        \node[state, accepting, right of=2] (3) {\(3\)};
        \draw
          (1) edge[below] node{\(a\)} (2)
          (2) edge[below] node{\(b\)} (3)
          (3) edge[bend right, above] node{\(b\)} (1);
      \end{tikzpicture}
    }
    \only<2> {
      \begin{tikzpicture}
        \tikzset{
          ->,
          >=stealth',
          node distance=2.25cm,
          every state/.style={
            thick,
            inner sep = 0pt,
            minimum size = 0.75cm,
          },
          initial text=$ $,
        }
        \node[state, initial] (1) {\(1\)};
        \node[state, above of=3] (1a) {\(1_a\)};
        \node[state, right of=1] (2b) {\(2_b\)};
        \node[state, above of=2b] (2) {\(2\)};
        \node[state, accepting, right of=2b] (3) {\(3\)};
        \node[state, accepting, above of=1] (3b) {\(3_b\)};
        \draw
          (1) edge[below] node{\(b\)} (2b)
          (2b) edge[below] node{\(a\)} (3)
          (3) edge[right] node{\(a\)} (1a)
          (1a) edge[above] node{\(b\)} (2)
          (2) edge[above] node{\(b\)} (3b)
          (3b) edge[left] node{\(b\)} (1);
      \end{tikzpicture}
    }
  \end{example}
\end{frame}

\begin{frame}{\myframetitle}
  \begin{block}{Compléxité en état de notre algorithme}
    Dans le pire cas, notre algorithme dupliquera tous les états de
    l'automate. 

    \pause[]
    \vphantom{}

    Sachant que chaque état peut être dupliqué étant de voir qu'il y a de
    symbole de l'alphabet plus un. 

    \pause[]
    \vphantom{}

    Alors, la complexité dans le pire cas de notre algorithme est \(n(\lvert
    \Sigma \rvert + 1)\).
  \end{block}
\end{frame}

% !chktex-file 8
% !chktex-file 41

\section{Conclusion}

\setframetitle{Conclusion}

\begin{frame}{\myframetitle}
  \begin{block}{Conclusion}
    On a donc que pour un alphabet binaire~:
    \begin{itemize}
      \item la concaténation de deux langages a une complexité de
        \(m2^n - 2^{n - 1}\),
      \item le renverser (à un état initial) d'un langage à une complexité de
        \((n + 1)\),
    \end{itemize}
  \end{block}
\end{frame}

\begin{frame}{\myframetitle}
  \begin{block}{Ouverture}
    On pourrait s'intéresser aux opérations~:
    \begin{itemize}
      \item<2-> celle de permutation (\textit{Twist}),
      \item<3-> celle de grignotage (\textit{Trognon}).
    \end{itemize}
  \end{block}
\end{frame}

\end{document}
